\documentclass[11pt]{article}
\usepackage[margin=25mm]{geometry}
\usepackage{fontspec}
\setmainfont{Latin Modern Roman}
\usepackage{amsmath,amssymb,mathtools}
\usepackage{unicode-math}
\setmathfont{Latin Modern Math}
\usepackage{xurl,hyperref}
\hypersetup{colorlinks=true,urlcolor=blue}
\setcounter{secnumdepth}{0}
\setlength{\parindent}{0pt}
\setlength{\parskip}{5pt}
\setlength{\emergencystretch}{3em}
\newcommand{\sref}[2]{\hyperlink{source:#1:#2}{[S#2]}}
\newcommand{\cataloguescope}[1]{\paragraph{Recorded target.}#1}
\newcommand{\R}{\mathbb R}
\newcommand{\C}{\mathbb C}
\newcommand{\N}{\mathbb N}
\newcommand{\Z}{\mathbb Z}
\newcommand{\Q}{\mathbb Q}
\newcommand{\F}{\mathbb F}
\newcommand{\eps}{\varepsilon}
\begin{document}
% UnsolvedMath ID 8100005; source catalogue code AMR-080-0005
\setcounter{section}{8100004}
\section{initial boundary value problems for integrable systems (IBVP)}\label{8100005}
{\small UnsolvedMath ID 8100005 · Mathematical Physics}
\subsection{Definitions and mathematical statement}

\paragraph{Shared notation.}$\mathbb N=\{1,2,\ldots\}$; the standard number fields have their usual meanings. All additional conventions are stated locally.


\paragraph{Definitions and notation.}On $x\ge0$, initial data and Dirichlet boundary data for NLS leave $u_x(0,t)$ as dependent data in the Fokas representation. A Robin condition is $u_x(0,t)+qu(0,t)=0$. A classical solution is twice continuously differentiable in $x$ and once in $t$ in the interior. Matching conditions concern compatibility of the initial and boundary traces at $(0,0)$. Finite-gap solutions are the quasiperiodic solutions supplied by periodic integrable spectral theory.

\paragraph{Catalogue wording retained for traceability.}
Determine the long-time behavior of focusing NLS $iu_t=u_{xx}+2|u|^2u$ on $x\ge0$ with given $u(x,0)=f(x)$ and $u(0,t)=e^{iwt}$, $w\in\mathbb R$. In particular, explain excitation of finite-gap modes as the driving frequency varies, analogous to the source’s periodically driven Toda chain $\ddot x_k=e^{x_{k-1}-x_k}-e^{x_k-x_{k+1}}$, $x_k(0)=\dot x_k(0)=0$ for $k\ge1$, $x_0(t)=2at+h(\gamma t)$, with $h$ $2\pi$-periodic, $\gamma>0$ and $a<a_{\mathrm{crit}}$. Use the Bikbaev--Tarasov method for defocusing NLS $iu_t=u_{xx}-2|u|^2u$ with Robin boundary conditions, including the singular extension and long-time behavior. For general appropriate smooth nonlinearities, Robin data and smooth matching initial data, prove that the half-line solution is classical in the interior; more generally determine smoothness for dispersive initial-boundary problems in any dimension.
\paragraph{Scope and source qualification.}The truncated ending and missing Toda driver hypotheses are restored from primary Problem 5. The requested admissible nonlinearities and regularity theory remain to be specified; no unconditional global existence theorem is inserted.
\subsection{Short English statement}
Determine driven half-line NLS asymptotics, solve the defocusing Robin problem, and prove interior smoothness for compatible dispersive boundary problems.
\subsection{Research attempt}
\paragraph{Model, date and investigated scope.}
OpenAI GPT-6 Astra (configured model \texttt{gpt-6-astra}), ultra reasoning; generated 11 October 2026. Constructed exact half-line focusing standing waves for w<=-1 and displayed two Neumann traces at the same frequency when w<-1. General driven asymptotics and all Robin/regularity/Toda questions remain unproved.

\paragraph{Sources and literature check.}
Read the supplied complete definition, catalogue wording, scope qualifications and source list for this canonical ID. The original-source citations are inherited as supplied references; no fresh primary-source or novelty search was performed for this attempt. No current open-status claim is inferred from this calculation.

\paragraph{Exact boundary-driven standing waves.}Use precisely the source's focusing sign \(iu_t=u_{xx}+2|u|^2u\). Seek \(u(x,t)=e^{iwt}\phi(x)\), real \(\phi\), with \(\phi(0)=1\) and \(\phi\to0\) at infinity. Then \(\phi''+2\phi^3+w\phi=0\). For \(\kappa>0\), direct differentiation gives
\[\phi(x)=\kappa\operatorname{sech}(\kappa(x+x_0)),\qquad w=-\kappa^2.
\]
The boundary condition requires \(\cosh(\kappa x_0)=\kappa\), so such profiles exist for \(\kappa\ge1\). For \(\kappa>1\) there are two choices \(x_0=\pm\kappa^{-1}\operatorname{arcosh}\kappa\); their Neumann traces are \(\phi'(0)=\mp\sqrt{\kappa^2-1}\). At \(\kappa=1\) the two choices coincide and the trace is zero. Taking initial data equal to either profile gives an exact classical solution with the specified unit-amplitude driver.
\paragraph{What this tests.}The same boundary frequency can coexist with different initial profiles and different dependent traces. A rule selecting the Neumann trace from frequency alone is therefore false without a condition on the incoming initial state or the asymptotic regime. These waves have exponential spatial decay, not compact support, and are not asserted to attract arbitrary data or to describe finite-gap excitation.
\paragraph{Remaining bundled scope.}For general given \(f\), long-time asymptotics and mode selection are still missing. The driven Toda analogy, the defocusing Robin problem and its singular extension, and interior smoothness for general compatible nonlinearities in any dimension are separate requirements not resolved by this explicit focusing family. The first gap here is control of perturbations of the profiles and the dependent boundary trace over long times. No defocusing threshold is imported into a focusing claim.

\paragraph{Verification and adversarial audit.}
Checked the source sign convention, sech derivative, amplitude-one boundary, two branches, w=-1 merger and the scope of the initial data. This was an LLM adversarial audit, not independent human verification; no theorem prover was used. Prompts A then B applied to this ID: Checked the source sign convention, sech derivative, amplitude-one boundary, two branches, w=-1 merger and the scope of the initial data. Scope remains limited: Control general initial-data evolution and its asymptotic boundary trace, including stability or selection of the explicit waves.

Coordinator checked the focusing standing-wave equation, boundary amplitude condition and the two Neumann traces. This additional check was model-assisted, not independent human verification.

\paragraph{Tools, process and contributions.}
Python 3 and local shell used to read supplied files and serialize this manuscript; mathematics derived symbolically in the text, with no numerical experiment or formal proof assistant. Contribution prompt A was followed by prompt B. GPT-6 Astra with ultra reasoning drafted the research and performed its stated self-audit; any additional coordinator contribution is identified in the verification and audit. A GPT-6-based Codex coordinator assembled the manuscript and checked repository integration; its exact variant and reasoning configuration were not exposed. No human mathematical contribution, independent human proof check or formal proof certification is supplied. Complete derivations and the consulted source references are included here; private reasoning transcripts are not published.

\paragraph{Final status and first remaining gap.}
\textbf{UNRESOLVED}. Control general initial-data evolution and its asymptotic boundary trace, including stability or selection of the explicit waves.

\section{Completion Estimate}
\noindent\textbf{COMPLETION: 0\%}
The calculation is an elementary benchmark or conditional reduction; it establishes no claimed new portion of the complete source target.

\subsection{Sources}
\begin{itemize}
\item[\hypertarget{source:8100005:1}{S1}] ULAM.AI and the UnsolvedMath Contributors, UnsolvedMath, record 8100005 (AMR-080-0005). Catalogue curation under CC BY 4.0; definitions and English exposition added locally.
\url{https://huggingface.co/datasets/ulamai/UnsolvedMath}.
\item[\hypertarget{source:8100005:2}{S2}] Percy Deift, Some Open Problems in Random Matrix Theory and the Theory of Integrable Systems. II, SIGMA 13 (2017), 016, Problem 5.
\url{https://arxiv.org/abs/1703.04931}.
\item[\hypertarget{source:8100005:3}{S3}] Primary full-text formulation and 2017 historical progress discussion.
\url{https://arxiv.org/html/1703.04931v1}.
\end{itemize}
\label{end:8100005}
\end{document}
